\section{Conclusion}
\label{sec:conclusion}

Nearly a decade after the original AE4 model, the biological conclusion remains compelling but the mechanistic explanation has changed. AE4 can carry a substantial stimulated chloride load only when the model includes an inward source of bicarbonate equivalents. NBC is therefore required not because it directly dominates fluid flow, but because it makes productive AE4 transport compatible with pH homeostasis. Once that wild type architecture is repaired, NKCC1 compensation largely removes the knockout phenotype. Neither equal cation routing nor the stoichiometric source classes suggested by recent molecular experiments recover the observed secretion loss under the inherited kinetics. A strong additional coupling to apical chloride conductance can reconstruct the magnitude, but is target selected and temporally incomplete.

The central conclusion is therefore not that one proposed AE4 stoichiometry is correct or incorrect. It is that AE4 control of secretion is a network property. Cation coordination, bicarbonate supply, intracellular pH, compensatory chloride loading and stimulus dependent regulation must be considered together. The 2018 model identified the right transporter. The present reconstruction identifies the balances that any complete explanation must satisfy.

\section*{Data and code availability}
All model source, frozen inputs, numerical outputs, validation records, figure scripts and the manuscript source are maintained in the accompanying GitHub repository. The prediction checkpoint for the Catalán mechanism class panel was committed before phenotype reveal. Exact commit identifiers are recorded in the repository reports.

\section*{Author contributions}
To be completed after the final author list is agreed.

\section*{Funding}
To be completed before submission.

\section*{Competing interests}
The authors declare no competing interests, subject to confirmation by all authors.

\bibliographystyle{plainnat}
\bibliography{references}
